% ============================================================
% CHAPTER 3 — The ADLC Method
% ============================================================
\chapter{The ADLC Method: How to Work in 0-Code}\index{ADLC}\index{Context Engineering}

\section*{What You Will Learn}
By the end of this chapter you will know:
\begin{itemize}
  \item The 7 phases of the Agent Development Life Cycle applied to application development
  \item How to write an effective \texttt{\_CONTEXT.md}
  \item Why Context Engineering is the key skill of the 0-code paradigm
  \item How to structure a project to maximize the quality of AI output
\end{itemize}

% ---------------------------------------------------------
\section{The Life Cycle: The 7 Phases of the ADLC}

The ADLC (Agent Development Life Cycle) is the structured method for developing with AI. It is not abstract theory --- it is the sequence of steps you will follow throughout the book.

\begin{center}
\begin{small}
\begin{tabularx}{\textwidth}{c l X X}
\toprule
\textbf{Phase} & \textbf{Name} & \textbf{What you do} & \textbf{Deliverable} \\
\midrule
\textbf{0} & Preparation & Define the problem and assess whether 0-code is suitable & Solution hypothesis \\
\textbf{1} & Framing & Establish boundaries, technologies, and constraints & Natural-language requirements \\
\textbf{2} & Definition & Write \texttt{\_CONTEXT.md} and any \texttt{SKILL.md} files & Context documents \\
\textbf{3} & Simulation & Validate the idea with a small case & Proof of value \\
\textbf{4} & Implementation & The AI generates code while you supervise & Working code \\
\textbf{5} & Release & Test, fix, and deploy & Production-ready application \\
\textbf{6} & Learning & Update the context with lessons learned & Improved context \\
\bottomrule
\end{tabularx}
\end{small}
\end{center}

\begin{notebox}
\textbf{Toward the professional framework.} In this chapter we work with a single \texttt{\_CONTEXT.md}. In Appendix~E you will see how the exact same principles scale to a professional framework with multiple modules, specialist skills, and infrastructure governance.
\end{notebox}

\subsection{How it applies in daily practice}

For a simple project, phases~0--3 may take only a few minutes. For a complex full-stack application, each phase becomes a meaningful checkpoint.

\textbf{Example --- a blog REST API}
\begin{itemize}
  \item \textbf{Phase 0}: ``I need a backend for a blog with posts, comments, and authentication.''
  \item \textbf{Phase 1}: ``I will use Node.js, Express, PostgreSQL, Prisma, and JWT.''
  \item \textbf{Phase 2}: Write the \texttt{\_CONTEXT.md} with architecture, structure, constraints, and commands.
  \item \textbf{Phase 3}: Ask the AI to generate only the \texttt{/health} endpoint.
  \item \textbf{Phase 4}: Expand to routes, services, middleware, and data model.
  \item \textbf{Phase 5}: Test, fix, and deploy.
  \item \textbf{Phase 6}: Update the context with what the AI got wrong.
\end{itemize}

\begin{quote}
Phase~6 is essential: the \texttt{\_CONTEXT.md} is a \textbf{living document}.
\end{quote}

% ---------------------------------------------------------
\subsection{When the AI goes off the rails}

Sooner or later the AI will enter an error loop: it fixes one issue, creates another, then breaks something that was previously working.

\begin{warnbox}
\textbf{Five recovery techniques when the AI gets stuck:}
\begin{enumerate}
  \item \textbf{Stop and reset the conversation}
  \item \textbf{Return to the last working version}
  \item \textbf{Break the task into smaller parts}
  \item \textbf{Add a more specific constraint to the context}
  \item \textbf{Switch model or tool if needed}
\end{enumerate}
\end{warnbox}

\begin{tipbox}
If the AI has failed twice on the same problem with the same approach, do not insist. Change the strategy.
\end{tipbox}

\subsection{Phase 6 in practice: how a constraint saves the project}

\begin{dangerbox}
\textbf{Disaster 1: passwords stored in plain text.} The AI generates a registration endpoint that stores user passwords unencrypted. The fix is not just changing the code once --- it is adding a permanent rule to the context:\\[4pt]
\texttt{SEC-01: Passwords MUST be hashed with bcrypt (cost 12). Never store passwords in plain text.}
\end{dangerbox}

\begin{dangerbox}
\textbf{Disaster 2: the AI changes stack mid-project.} You are using ESModules and the AI suddenly generates CommonJS. The preventing rule is explicit:\\[4pt]
\texttt{Language: JavaScript ES2022+ with ESModules. DO NOT use require() or module.exports.}
\end{dangerbox}

\begin{dangerbox}
\textbf{Disaster 3: phantom dependencies.} The AI installs new libraries without approval. Prevent it with:\\[4pt]
\texttt{Dependencies: ONLY those listed in the \_CONTEXT.md. Do not add new dependencies without explicit approval.}
\end{dangerbox}

% ---------------------------------------------------------
\section{Context Engineering: The Art of Communicating with AI}

Context Engineering is the foundational discipline of 0-code development. It is the difference between mediocre code and professional-quality output.

\subsection{The amnesia problem}

Large language models are \textbf{stateless}. When you start a new chat, they do not remember:
\begin{itemize}
  \item the architectural decisions from yesterday,
  \item the naming conventions of your project,
  \item the bugs you already fixed together,
  \item the exact stack you selected.
\end{itemize}

If you do not provide that information, the model will invent its own assumptions.

\subsection{The solution: the \texttt{\_CONTEXT.md} file}

The \texttt{\_CONTEXT.md} file is the \textbf{contract} between you and the AI. It tells the agent what the project is, how it is structured, what conventions it must follow, and what it must never do.

\begin{quote}
Think of \texttt{\_CONTEXT.md} as the briefing document you would hand to a new developer on their first day.
\end{quote}

% ---------------------------------------------------------
\section{Anatomy of a \texttt{\_CONTEXT.md}}

A good context file usually contains five sections.

\subsection{1. Identity and Purpose}
\begin{lstlisting}[language={}]
# Project: BlogAPI

You are working on a REST API backend for a personal blog.
Technologies: Node.js 20, Express.js, PostgreSQL 16, Prisma ORM.
Environment: local development on Windows/macOS.
\end{lstlisting}

\subsection{2. Architecture and Structure}
\begin{lstlisting}[language={}]
## Project Structure

src/
  routes/          -> Endpoint definitions
  controllers/     -> Request/response logic
  services/        -> Business logic
  middleware/      -> Auth and error handling
  utils/           -> Helper functions
prisma/
  schema.prisma    -> Database schema
tests/
  unit/
  integration/
\end{lstlisting}

\subsection{3. Conventions and Standards}
\begin{lstlisting}[language={}]
## Conventions
- Naming: camelCase for variables and functions
- API responses: { success: boolean, data: T, error?: string }
- Validation: use Zod for input validation
- Async: use async/await, never callbacks
- Imports: use ESModules, never require()
\end{lstlisting}

\subsection{4. Constraints and Prohibitions}
\begin{lstlisting}[language={}]
## Constraints (MANDATORY)
- NEVER use raw SQL queries. ALWAYS use Prisma.
- NEVER store passwords in plain text.
- NEVER expose stack traces in production errors.
- NEVER install dependencies without a documented reason.
\end{lstlisting}

\subsection{5. Commands and Workflow}
\begin{lstlisting}[language={}]
## Commands
- Start the server: npm run dev
- Run tests: npm test
- Migrate database: npx prisma migrate dev
- Generate Prisma client: npx prisma generate
\end{lstlisting}

% ---------------------------------------------------------
\section{Your First Complete \texttt{\_CONTEXT.md}}

\begin{praticabox}[title={Hands-on --- Create your first context file}]
Create a \texttt{\_CONTEXT.md} file in the root of a new project:
\end{praticabox}

\begin{lstlisting}[language={},title={\texttt{\_CONTEXT.md} for TaskMaster CLI}]
# Project: TaskMaster CLI

## Purpose
A command-line application for managing personal tasks.
The user can add, list, complete, and delete tasks.
Data is saved in a local JSON file.

## Technologies
- Language: Python 3.11+
- Libraries: standard modules only (argparse, json, os, datetime)
- No external dependencies allowed

## Structure
taskmaster/
├── _CONTEXT.md
├── main.py
├── task_manager.py
├── storage.py
├── tests/
│   ├── test_task_manager.py
│   └── test_storage.py
└── data/
    └── tasks.json

## Conventions
- Naming: snake_case for everything
- Docstrings on all public functions
- Type hints on all function signatures
- Clear CLI output for the user

## Constraints
- DO NOT use external libraries.
- DO NOT use global variables.
- DO NOT ignore errors.
- tasks.json must be created automatically.
\end{lstlisting}

\begin{checkpointbox}
This is enough context for the AI to generate a working, tested CLI application.
\end{checkpointbox}

% ---------------------------------------------------------
\section{The 7 Golden Rules of Context Engineering}

\begin{enumerate}
  \item \textbf{Be imperative, not descriptive}
  \item \textbf{Prohibitions matter more than permissions}
  \item \textbf{Show the structure, do not vaguely describe it}
  \item \textbf{Specify the exact commands}
  \item \textbf{Use one \texttt{\_CONTEXT.md} per project}
  \item \textbf{Update it after every important error}
  \item \textbf{Less is more, as long as it is precise}
\end{enumerate}

% ---------------------------------------------------------
\section{MCP and Agent Skills: A Preview}

\subsection{Model Context Protocol (MCP)}
When the AI needs to interact with external systems --- databases, APIs, GitHub, or other tools --- it uses the \textbf{Model Context Protocol}. We will introduce it concretely in Chapter~7.

\subsection{Agent Skills (\texttt{SKILL.md})}
As projects become more complex, you can provide specialist expertise with \texttt{SKILL.md} files. Later chapters will use these for React, Flutter, API design, database work, and security.

% ---------------------------------------------------------
\section{Managing Long Projects: \texttt{PROGRESS.md}}

For projects that span multiple sessions, you need a persistent memory file.

\begin{lstlisting}[language={}]
Update PROGRESS.md with:
1. What was implemented in this session
2. What architectural decisions were made
3. What problems were encountered and how they were solved
4. What remains to be done
\end{lstlisting}

\begin{notebox}
With Claude Code you can also use a file such as \texttt{claude-progress.txt}. The principle is the same: persistent context between sessions.
\end{notebox}

% ---------------------------------------------------------
\section{The Professional Workflow: From Idea to Implementation}

Real projects require one more step: use the AI to decompose the idea into implementable work units before writing the code.

\begin{figure}[H]
\centering
\includegraphics[width=0.5\textwidth]{11-workflow-adlc}
\caption{Professional ADLC workflow: from idea to implementation in six phases}
\label{fig:workflow-adlc-en}
\end{figure}

\subsection{Phase 1: From idea to specifications}
Ask the AI to generate a PRD with functional requirements, non-functional requirements, technical constraints, and acceptance criteria.

\subsection{Phase 2: From requirements to use cases}
Derive concrete use cases from the PRD: actor, precondition, main flow, alternative flows, and postcondition.

\subsection{Phase 3: Design and technology choices}
The AI proposes options, but the human decides. Capture the result in \texttt{docs/DESIGN.md} using clear ADRs.

\subsection{Phase 4: From use cases to epics}
Group related work into epics such as setup, authentication, CRUD, frontend, mobile, testing, and deployment.

\subsection{Phase 5: From epics to tasks}
Break each epic into atomic tasks with clear acceptance criteria, involved files, SEC/PERF constraints, and an estimate.

\subsection{Phase 6: Task-driven implementation}
Once the tasks are defined, implementation becomes systematic and traceable.

\begin{center}
\begin{tabularx}{\textwidth}{X X}
\toprule
\textbf{Without decomposition} & \textbf{With decomposition} \\
\midrule
``Build the BookShelf app'' & ``Implement T-003.1'' \\
The AI invents the architecture & The architecture is documented first \\
If it goes off track, start over & If it goes off track, resume from the task \\
No traceability & Every decision is recorded \\
Hard to estimate work & Story points per task \\
\bottomrule
\end{tabularx}
\end{center}

\begin{checkpointbox}
The workflow Idea $\rightarrow$ Specifications $\rightarrow$ Use Cases $\rightarrow$ Design $\rightarrow$ Epics $\rightarrow$ Tasks $\rightarrow$ Implementation is what turns AI experimentation into software engineering.
\end{checkpointbox}

% ---------------------------------------------------------
\section*{Summary}

\begin{center}
\begin{tabularx}{\textwidth}{l X}
\toprule
\textbf{Concept} & \textbf{What you learned} \\
\midrule
ADLC & The 7 phases of the 0-code life cycle \\
\texttt{\_CONTEXT.md} & How to write the contract with the AI \\
Context Engineering & The rules for high-quality prompts and constraints \\
Project structure & The five essential parts of a well-written context \\
MCP / Skills & A preview of advanced agent tooling \\
\texttt{PROGRESS.md} & Persistent memory between sessions \\
Professional workflow & How to decompose an idea into implementable tasks \\
\bottomrule
\end{tabularx}
\end{center}
